\chapter{Justificación}\label{cap2}

Los ascomicetos coprófilos representan una reserva importante de antibióticos y
otros metabolitos secundarios bioactivos, útiles para el desarrollo de nuevos
tratamientos farmacéuticos \cite{Bills2013}. Estas propiedades cobran relevancia debido al
aumento de las bacterias multirresistentes y la limitada eficacia de antibióticos
actuales. Aunque antibióticos como la penicilina provienen de hongos, desde la
década de 1990, barreras técnicas y regulaciones han limitado el descubrimiento
de nuevos compuestos fúngicos \cite{Correia2023}.

A pesar del conocido potencial biotecnológico de estos hongos, rara vez han sido
aislados para ser cultivados de forma controlada con fines experimentales, a fin de
comprender su comportamiento y, posteriormente, fomentar la producción de
biomasa y metabolitos \cite{MarinFelix}. En la naturaleza, el ecosistema óptimo para la
proliferación de estos hongos son los matorrales xerófilos, los cuales son
vulnerables al deterioro por el cambio de uso del suelo \cite{Sanchez2004}.

Actualmente, en el área de Bioquímica de la Escuela Nacional de Ciencias
Biológicas (ENCB) del Instituto Politécnico Nacional (IPN), micólogos estudian
hongos coprófilos en cajas Petri. Sin embargo, al escalar el cultivo en biorreactores,
se presentan diferentes problemáticas, puesto que se carece de equipo
especializado que considere las características de estos hongos, aseguren su
crecimiento uniforme y sin obstrucciones, y faciliten la limpieza del equipo y
recolección de la biomasa generada.

En este contexto, se identifica la necesidad de generar una estación de pruebas
controladas para la investigación de hongos ascomicetos coprófilos. Para ello, se
propone diseñar y construir un biorreactor de tipo airlift, uno de los preferidos para
el cultivo de hongos filamentosos, debido a su eficiente transferencia volumétrica
de oxígeno y a que evita el estrés por cizallamiento generado por la agitación
mecánica presente en otro tipo de biorreactores \cite{Cerrone2025}.

El desarrollo de este proyecto representará una herramienta práctica en el campo
de la biotecnología fúngica, puesto que brindará un ambiente cerrado y controlado,
asegurando la reproducibilidad de los experimentos, lo que facilitará la
experimentación científica en este tipo de hongos.
